\documentclass[11pt, a4paper, twoside, titlepage]{book}

% Paquetes usados

\usepackage[a4paper, top=4cm, bottom=4cm, left=3.3cm, right=3.2cm]{geometry}
%\usepackage[dvips]{graphicx} %%latex -> Usamos  .eps
\usepackage{graphicx} %%pdflatex -> Usamos .jpg/.png
\usepackage{fancyhdr}
\usepackage{subfigure}
\usepackage{booktabs}
\usepackage{url}
\usepackage{listings}
\usepackage{enumitem}
\usepackage{float}
\usepackage{eurosym}
\usepackage{colortbl}

\usepackage[draft]{todonotes}  %% CVR  draft/disable (borrador/definitivo)
%\todo[inline]{comment}  %% CVR, para poner un comentario, [inline] es opcional
%\listoftodos %% CVR, para mostrar todos los comentarios
%\missingfigure{Add my picture here.}


%\usepackage[latin1]{inputenc} %% CV
\usepackage[utf8]{inputenc}  %% CVR
%\usepackage[spanish]{babel} %% CVR
\usepackage[spanish,activeacute]{babel} %% CVR

\usepackage{color}   %May be necessary if you want to color links
\usepackage{hyperref}
\hypersetup{
    colorlinks=true, %set true if you want colored links
    linktoc=all,     %set to all if you want both sections and subsections linked
    linkcolor=blue,  %choose some color if you want links to stand out
}

\newlist{cu}{enumerate}{1}
\setlist[cu,1]{
	label={\textbf{CU-\arabic*~}},
	leftmargin=*
	%align=left,
	%labelsep=10mm,
	%itemindent=4em
}

\definecolor{pblue}{rgb}{0.13,0.13,1}
\definecolor{pgreen}{rgb}{0,0.5,0}
\definecolor{pred}{rgb}{0.9,0,0}
\definecolor{pgrey}{rgb}{0.46,0.45,0.48}

\lstset{
	language=Java,
	showspaces=false,
	showtabs=false,
	breaklines=true,
	showstringspaces=false,
	breakatwhitespace=true,
	frame=single,
	tabsize=2,
	commentstyle=\color{pgreen},
	keywordstyle=\color{pblue},
	stringstyle=\color{pred},
	basicstyle=\ttfamily,
	moredelim=[is][\color{pgrey}]{$$}{$$},
	basicstyle=\footnotesize
}


% Aumentar la separación entre líneas
\renewcommand{\baselinestretch}{1.5}


% Para eliminar la cabecera de las páginas vacías al final de los
% capítulos
\makeatletter
\def\cleardoublepage{\clearpage\if@twoside \ifodd\c@page\else
  \hbox{}
    \thispagestyle{empty}
      \newpage
        \if@twocolumn\hbox{}\newpage\fi\fi\fi}
\makeatother
%%


\begin{document}


\pagestyle{empty}

\begin{titlepage}
\begin{center}

\includegraphics[scale=0.4]{figures/udc}


\Large{FACULTADE DE INFORMÁTICA}

\emph{\large{Departamento de Enxeñería de Computadores}}

Proxecto de Fin de Carreira de Enxeñería Informática

\vspace*{3cm}

\textbf{\LARGE{Caronte: plataforma para o guiado no interior de museos}}

\end{center}

\vspace*{4cm}

\begin{flushright}
\large{
\emph{Autor:} Daniel Pan Farto\\
\emph{Titor: } Carlos Vázquez Regueiro\\
Culleredo, \today}
\end{flushright}

\end{titlepage}

\cleardoublepage

% Numeración romana de páginas
\pagenumbering{roman}


% Especificación del pfc
\chapter*{Especificación}{}


\begin{tabular}{lp{9cm}}
\emph{Título}: & Caronte: plataforma para o guiado no interior de museos \\
& \\
\emph{Clase}: & Proxecto clásico de enxeñería \\
& \\
\emph{Autor}: &  Daniel Pan Farto \\
& \\
\emph{Director}: & Carlos Vázquez Regueiro \\
& \\
\emph{Tribunal}: & \\
& \\
& \\
& \\
& \\
& \\
& \\
\emph{Data de lectura}: & \\
& \\
\emph{Calificación}: & \\
\end{tabular}


% Dedicatoria y agradecimientos
\include{agradecimientos}



% Resumen y palabras clave
\chapter*{Resumo}

Poucos elementos tecnolóxicos entraron máis rápido na vida cotiá da xente que os teléfonos móbiles intelixentes, conseguindo facerse indispensábeis nun curto espazo de tempo. Son moitas as aplicacións que teñen no día a día das persoas, chegando a substituír por completo os elementos máis tradicionais cos que se levaban a cabo anteriormente. Entre estas aplicacións encóntrase a xeolocalización. Grazas a tecnoloxías coma o GPS (Global Positioning System) é posíbel situar un dispositivo en calquera punto do planeta cunha pequena marxe de erro, permitindo non só a localización dun elemento senón a indicación de rutas entre dous puntos.

Un dos principais problemas dos sistemas coma o GPS é a imposibilidade de seren utilizados en interiores. A localización en interiores presenta outro problema e é a precisión que se require para un correcto funcionamento, pois o erro asumíbel do GPS en exterior non sería viábel para os espazos interiores.

Nos últimos anos están aparecendo diversos sistemas de localización en interiores mediante a utilización de teléfonos móbiles intelixentes, sendo algúns deles moi precisos no seu posicionamento. Poderíanse describir como “sistemas GPS de interiores”. Non usan os sinais fornecidos polos satélites GPS, senón que utilizan outro tipo de sinais accesíbeis dende o teléfono móbil.

Grazas a esta posibilidade de localización en interiores é posíbel a navegación, é dicir, a creación dunha ruta entre distintas posicións que indique a unha persoa o movemento a realizar.

A cantidade de información que pode haber ao redor dunha obra de arte pode ser abrumadora: dende as técnicas utilizadas para a creación dunha escultura, como o contexto histórico do artista ou o movemento artístico ao que pertence unha pintura. Cos métodos tradicionais, non sería posíbel mostrar toda a información dispoñíbel sobre unha obra nos museos. As guías de papel teñen espazo limitado e só poden mostrar imaxes e texto; e as guías con audio requiren a introdución de códigos para a súa utilización e non proporcionan máis información.

O tempo dispoñíbel á hora de percorrer un museo de gran tamaño adoita ser un problema para os visitantes. Non dispoñer do tempo desexado obriga á escolla das obras máis interesantes segundo a opinión do visitante e non sempre é doado realizala. Hai veces nas que o visitante desexa é admirar as obras de certos artistas ou movementos artísticos, o cal pode ser complicado se non posúe moita información.

\section*{Palabras clave}

Android, Localización en interiores, Situm Technologies, Guiado en museos.



% Definimos el encabezado y pie de página
\pagestyle{fancy}
\renewcommand{\chaptermark}[1]{%
\markboth{\thechapter.\ #1}{}}
\renewcommand{\sectionmark}[1]{%
\markright{\thesection.\ #1}{}}
\fancyhead{}
\fancyhead[LE,RO]{\thepage}
\fancyhead[LO,ER]{\rightmark}
\fancyfoot{}


% Índice de capítulos, secciones y subsecciones
\cleardoublepage
\pdfbookmark[0]{\contentsname}{toc}%
\tableofcontents

% Índice de figuras
\cleardoublepage
\phantomsection \label{listoffigures}
\addcontentsline{toc}{chapter}{\listfigurename}
\listoffigures


% Redefinición del nombre que encabeza las tablas. Por defecto es cuadro.
\renewcommand\tablename{Tabla}
\renewcommand\listtablename{Índice de tablas}

% Índice de tablas 
\cleardoublepage
\phantomsection \label{listoftables}
\addcontentsline{toc}{chapter}{\listtablename}
\listoftables
\cleardoublepage


% Numeración normal de páginas
\pagenumbering{arabic}


% Definimos el encabezado y pie de página
\fancyhead{}
\fancyhead[LE,RO]{\thepage}
\fancyhead[LO]{\rightmark}
\fancyhead[ER]{\leftmark}


% Capítulos

\include{introduccion}
\include{estado_del_arte}
\chapter{Fundamentos tecnolóxicos}
Neste capítulo farase unha relación das tecnoloxías hardware e software nas que se basea o proxecto e outras que están en clara conexión. Tamén se tratará todo aquel software utilizado para a creación deste proxecto.


\section{Recursos hardware}
Para a realización deste proxecto empregáronse certos recursos hardware que pasaremos a detallar a continuación:


\subsection{Dispositivos (teléfonos) móbiles intelixentes}
Un dispositivo (teléfono) móbil intelixente é un ordenador de tamaño reducido con capacidades de comunicación pola rede telefónica a través de voz e datos. Están deseñados para realizar múltiples tarefas a través de software instalado no mesmo, xa sexa polo usuario ou fabricante, que se apoia nun sistema operativo que permite a utilización do hardware incluído no teléfono.

Os primeiros dispositivos foron creados a finais do século pasado mais non comezaron a estenderse ata esta década. Na actualidade é tan habitual o uso destes dispositivos que xa se contan por miles de millóns en todo o planeta.

Debido ao seu uso tan estendido é a opción escollida para ser a base do proxecto.

\subsubsection{One Plus 3T}

O dispositivo utilizado para as probas da aplicación Android foi o terminal One Plus 3T, coa versión 8.0 do sistema Android, un terminal de gama media que permitirá facerse unha idea do desempeño xeral da aplicación neste tipo de dispositivos.

As súas características pódense observar na táboa \ref{tab:tabCaracteristicasOnePlus}

\begin{table} [tbh]
	\centering
	\footnotesize
	\begin{tabular}{|l|p{10cm}|}
		\hline 
		& \textbf{Características One Plus 3T} \\ 
		\hline 
		\textbf{Sistema Operativo} & Android 8.0 Oreo \\ 
		\hline 
		\textbf{Procesador} & Qualcomm Snapdragon 821 Quad Core, Kryo \\ 
		\hline 
		\textbf{Dimensións} & 152,7 x 74,7 x 7,35 mm \\ 
		\hline 
		\textbf{Peso} & 158 g \\ 
		\hline 
		\textbf{Pantalla} & 5,5 pulgadas cunha resolución 1080p Full HD (1920 x 1080 píxeles) \\ 
		\hline 
		\textbf{Memoria} & 6 GB \\ 
		\hline
		\textbf{Almacenamento} & 64 GB \\ 
		\hline 
	\end{tabular}
	\caption{Características One Plus 3T}
	\label{tab:tabCaracteristicasOnePlus}
\end{table}


\section{Recursos software específicos do proxecto}
Estes recursos foron básicos na creación do sistema, xa sexa para a elaboración do código ou como plataforma para executar o noso código. A continuación pásase a explicar brevemente a súa función:

\subsection{SDK Situm}
Situm é unha empresa galega especializada no posicionamento en interiores. O seu sistema pódese implantar en calquera edificio cunha infraestrutura mínima, xa que incluso pode funcionar sen instalación de ningún tipo de dispositivos. Grazas a iso é posíbel un despregue inmediato no que só sería preciso a calibración do edificio que se pode levar a cabo en minutos. Situm proporciona un SDK (Software Development Kit - ou en galego, Kit de desenvolvemento de software) que permite utilizar o seu sistema tanto en Android coma iOS. Dispoñen dunha API pública na que poder revisar todos os elementos incluídos no SDK facilitando así o seu uso.


\subsection{Apache Tomcat}
Apache Tomcat é un contedor de servlets Java desenvolvido pola Apache Software Foundation. Implementa varias especificacións da Java Enterprise Edition e proporciona un entorno para servidores web HTTP no cal se pode executar código Java. Foi a opción escollida para o despregue dos servizos.


\subsection{Amazon Web Services}
Amazon Web Services (AWS) é unha plataforma de servizos na nube que ofrece gran potencia de cómputo, almacenamento para bases de datos, entrega de contido, entre outras funcionalidades aportando gran flexibilidade, un entorno escalábel e moi fiábel. Esta plataforma supón un gasto que depende da capacidade de cómputo requirida así coma o uso de rede. Utilizouse esta plataforma para lanzar o noso servidor e almacenar a base de datos.


\subsection{Eclipse}
Eclipse é un conxunto de ferramentas, proxectos e grupos de traballo en código aberto. Entre as ferramentas das que se compón, unha das máis utilizadas é o seu IDE (Integrated Development Environment - ou en galego, Entorno de desenvolvemento integrado), dispoñíbel para distintos sistemas operativos e amplamente utilizado para a programación en Java. Este foi o entorno que se utilizou para a elaboración dos servizos web.


\subsection{Java}
A plataforma Java é un entorno de computación creado por Sun Microsystems capaz de executar aplicacións desenvolvidas na linguaxe de programación Java, principalmente. A encargada de executar as aplicacións é a súa máquina virtual xunto cunha serie de bibliotecas estándar. Foi utilizada conxuntamente coa linguaxe de programación do mesmo nome para a creación dos servizos web.


\subsection{Android Studio}
É o IDE oficial para a programación en Android. Dende a súa saída substituíu ao Eclipse que era o anterior. Ao estar especificamente deseñado para o desenvolvemento de aplicacións en Android fai que a experiencia de programación para este sistema sexa moito máis rápida e cómoda. Foi o entorno utilizado para a programación da aplicación de Android.


\subsection{Postman}
Postman é un IDE destinado ao desenvolvemento de APIs. Permite a creación de peticións a APIs e a elaboración de tests de validación do seu comportamento entre outras características. Utilizouse para as probas da parte servidora.


\subsection{SQuirreL}
SQuirreL é un cliente de SQL que utiliza Java para se executar. A súa interface de acceso á base de datos é gráfica, o que permite ver a estrutura da mesma, así como navegar a través dos seus esquemas e táboas. A súa funcionalidade pode estenderse a través de plugins. Foi a aplicación utilizada para a xestión da base de datos.


\section{Recursos de xestión do proxecto}
A continuación faise unha relación dos recursos utilizados na xestión do proxecto ou na elaboración da memoria:

\subsection{Git}
Git é un sistema de control de versións que permite a coordinación de distintas persoas que realizan cambios sobre o mesmo traballo. Está orientado á rapidez, manter a integridade dos datos e soportar modelos de traballo distribuídos e non lineais. Como servidor para o control de versións utilizouse GitHub.


\subsection{LaTeX e TeXstudio}
LaTeX é un sistema de composición de textos orientado á creación de documentos cunha alta calidade tipográfica. Polas súas características e posibilidades, úsase especialmente na xeración de artigos e libros científicos que inclúen expresións matemáticas. Editoriais científicas de primeira liña utilizan este sitema debido á súa calidade.
TeXstudio é un IDE para a creación de documentos en LaTeX, cun deseño simple e agradábel, que integra numerosas ferramentas para a creación de documentos científicos. É de código aberto. Utilizouse este entorno para a elaboración da memoria.


\subsection{Inkscape}
É un editor de gráficos vectoriais, libre e de código aberto. Permite a conversión de ficheiros de imaxes clásicos en documentos de gráficos vectoriais. Utilizouse para a conversión das imaxes utilizadas na memoria.


\subsection{StarUML}
StarUML é un proxecto de código aberto que permite a creación de calquera tipo de diagrama UML. Utilizouse para a elaboración dos diagramas incluídos na memoria.

\include{fund_tecnologicos}
\include{analisis}
\include{metodologia}
\include{gestion}
\include{arquitectura}
\include{pruebas}
\include{conclusiones}


% Apéndices
\appendix

\include{apen1}


\include{apen_situm}

\include{apen_contenidoDVD}



% Definimos el encabezado y pie de página
\fancyhead{}
\fancyhead[LE,RO]{\thepage}
\fancyhead[LO,ER]{\rightmark}

% Glosario de términos
%\printglossary


\cleardoublepage
\addcontentsline{toc}{chapter}{Bibliografía}
\cleardoublepage

% Bibliografía
\bibliographystyle{pfc-fic}
\bibliography{biblio}

\end{document}
